\chapter*{Resumen}
\addcontentsline{toc}{chapter}{Resumen}

Este Trabajo Fin de Máster estudia dos preguntas que conviene separar. La primera es si un sistema
de aprendizaje automático formado por cinco agentes especializados puede ordenar acciones de mejor a
peor y conservar esa capacidad fuera de muestra. La segunda es si esa ordenación puede cobrarse
frente al índice ---un objetivo de suma cero antes de costes, en el que la mayoría de los fondos
profesionales fracasa--- y cuánto depende esa respuesta de las reglas que convierten el orden en una
cartera. El sistema trabaja con un universo histórico del S\&P 500, fundamentales fechados por su
publicación y cohortes cerradas. Tres Model Studies encadenados seleccionan el modelo exclusivamente
por Rank-IC; después, con la señal congelada, un Portfolio Study compara \NumCarteras{}
configuraciones mediante Information Ratio.

En 2015--2024, el ganador alcanza un Rank-IC medio de \RankICSeleccion{} sobre
\CohortesSeleccion{} cohortes, con \FraccionICPositiva{} positivas, $t$ de Newey--West
\TNeweyWest{} y $p=\PValorPermutacion{}$ en \NumPermutaciones{} permutaciones. La era reservada
2025--2026 mantiene Rank-IC positivo (\RankICReservado{}), pero solo contiene
\CohortesReservadas{} cohortes. Además, el agente \texttt{risk} supera ligeramente al meta-agente,
por lo que la evidencia favorece la capacidad de ordenación, pero no demuestra que la arquitectura
multiagente sea necesaria.

La respuesta a la segunda pregunta resulta depender de la cartera hasta cambiar de signo. En
selección, el Information Ratio pasa de \IRModelo{} a \IRCartera{} y la rotación baja de
\TurnoverModelo{} a \TurnoverCartera{} veces al año. En la era reservada, la misma señal congelada
pasa de un exceso de \ExcesoModeloReservado{} con la cartera original a \ExcesoCarteraReservado{}
con la optimizada: de destruir valor a igualar al índice. Ese cambio de signo es alentador, pero no
confirma que el sistema bata al mercado: la cartera es la mejor de \NumCarteras{} alternativas y la
confirmación cubre solo \AnosCarteraReservada{} años. La contribución principal es, por tanto,
metodológica: una señal predictiva y su implementación económica deben evaluarse por separado y toda
conclusión práctica debe declarar la cartera con la que se midió.

\noindent\textbf{Palabras clave:} aprendizaje automático, selección de acciones, Rank-IC,
\textit{point-in-time}, construcción de cartera, validación fuera de muestra.

\clearpage
\chapter*{Abstract}
\addcontentsline{toc}{chapter}{Abstract}

This Master's Thesis addresses two questions that should not be conflated. First, can a machine
learning system built from five specialised agents rank equities and preserve that ability out of
sample? Second, can that ranking actually be monetised against the index ---a zero-sum objective
before costs, at which most professional funds fail--- and how strongly does the answer depend on the
rules that translate the ranking into a portfolio? The system uses a historical S\&P 500 universe,
filing-dated fundamentals, and closed cohorts. Three chained Model Studies select the predictive
model exclusively through Rank-IC. Once the signal is frozen, a Portfolio Study compares
\NumCarteras{} portfolio configurations using the Information Ratio.

Over 2015--2024, the selected model obtains a mean Rank-IC of \RankICSeleccion{} across
\CohortesSeleccion{} cohorts, with a Newey--West statistic of \TNeweyWest{} and a permutation
$p$-value of \PValorPermutacion{}. The reserved 2025--2026 period remains positive
(\RankICReservado{}), but contains only \CohortesReservadas{} closed cohorts. The standalone
\texttt{risk} specialist also ranks slightly better than the meta-agent, so the evidence supports
predictive ranking ability without establishing that the multi-agent architecture is necessary.

The answer to the second question turns out to depend on the portfolio to the point of changing
sign. In the selection window, the Information Ratio rises from \IRModelo{} to \IRCartera{} while
annual turnover falls from \TurnoverModelo{} to \TurnoverCartera{}. In the reserved period, the same
frozen signal moves from \ExcesoModeloReservado{} excess return under the original portfolio to
\ExcesoCarteraReservado{} under the selected one: from destroying value to matching the index. The
sign reversal is encouraging but does not establish that the system beats the market: the portfolio
is the best of \NumCarteras{} alternatives and the reserved evaluation spans only
\AnosCarteraReservada{} years. The main contribution is therefore methodological: predictive ranking
and economic implementation must be evaluated separately, and any practical claim must state the
portfolio through which the signal was measured.

\noindent\textbf{Keywords:} machine learning, equity selection, Rank-IC, point-in-time data,
portfolio construction, out-of-sample validation.
